\documentclass[uplatex,dvipdfmx,a4paper,11pt]{jsarticle}

\usepackage[top=25mm,bottom=25mm,left=25mm,right=25mm]{geometry}

\usepackage{enumitem}

\setlength{\parindent}{0pt}
\setlength{\parskip}{0.6em}

\usepackage{titlesec}
\titleformat{\section}
  {\normalfont\Large\bfseries}{\thesection.}{0.6em}{}
\titleformat{\subsection}
  {\normalfont\large\bfseries}{\thesubsection}{0.6em}{}
\titlespacing*{\section}{0pt}{1.2em}{0.6em}
\titlespacing*{\subsection}{0pt}{1.0em}{0.4em}

\renewcommand{\labelitemi}{\textbullet}
\renewcommand{\labelitemii}{\textopenbullet}
\setlist[itemize]{leftmargin=1.6em,itemsep=0.1em,topsep=0.2em}

\providecommand{\textopenbullet}{\ensuremath{\circ}}


\begin{document}

\section{Display（LEDマトリクス制御）変更履歴}

中継機のArduinoに搭載されているLEDMatrix制御を行うDisplayにおける計画書との変更点について説明する．まず，今回変更をした箇所以下の通りである．

\begin{itemize}
\item 描画APIの変更
\item 関数の追加
\item モジュール構成の変更
\item 静的IPの設定
\end{itemize}

まず，LEDを発光させる関数として，\texttt{matrix.renderBitmap()}から\texttt{matrix.loadFrame()}に変更した．この変更理由は，\texttt{matrix.renderBitmap()} は\texttt{uint8\_t[8][12]} 形式のビットマップ配列を入力として，ライブラリ内部で表示用フレーム形式に変換する必要があるのに対し，
\texttt{matrix.loadFrame()}事前にビットパックを行い，ライブラリ内部のフレーム形式である\texttt{uint32\_t[3]}形式で直接受け取るため，表示するデータの変換処理を明示的に制御できるためである．\par

次に，演奏終了に対応しLEDを消灯させる関数\texttt{clearDisplay()}と，\texttt{matrix.begin()}をDisplay.cpp側で呼ぶための初期化関数\texttt{diaplaySetup()}を追加した．
\texttt{clearDisplay()}の追加理由は，通信プロトコルにて演奏を終了させるヘッダ'E'が定義されているが，Display側の関数一覧には対応する処理への言及がされておらず，演奏終了後にLEDを消灯させる必要があったためである．
\texttt{diaplaySetup()}の追加理由は， \texttt{Arduino\_LED\_Matrix}の表示バッファはヘッダ内でstatic定義されており，includeする翻訳単位ごとに別実体となるため，
描画を行う\texttt{matrix.loadFrame()}と同一ファイルから\texttt{matrix.begin()}を呼び出す必要があったためである．\par
次に，\texttt{checkUDP()}をDisplay.inoからDisplay.cppに移動した．この変更理由は，\texttt{checkUDP()}内で\texttt{currentBPM}を参照する必要があったためである．\par
次に，\texttt{WiFi.config(localIP, gateway, subnet)} を\texttt{WiFi.begin()} の前に呼び出す処理を追加した．この変更理由は，
計画書には\texttt{localIP}が定数として定義されているが，Display側で静的IPを設定する具体的な手順が記載されていなかったためある．





 

\end{document}
